\chapter{Agents}
\label{ch:agents}

Agents are LLM-powered entities. They have a provider, a system prompt, optional tools, and optional memory. Agents are the primary way to interact with language models in \haira{}.

\section{Agent Declaration}

\begin{lstlisting}
agent AgentName {
    provider: provider_name
    system: "System prompt describing the agent's behavior."
    tools: [tool1, tool2]
    temperature: 0.5
    timeout: 60
    memory: conversation(max_turns: 50)
}
\end{lstlisting}

\section{Agent Fields}

\begin{table}[h]
\centering
\begin{tabular}{llll}
\toprule
\textbf{Field} & \textbf{Type} & \textbf{Required} & \textbf{Description} \\
\midrule
\code{provider} & identifier & Yes & Provider to use \\
\code{system} & \type{string} & Yes & System prompt \\
\code{tools} & \type{[tool]} & No & List of tools the agent can call \\
\code{temperature} & \type{float} & No & Sampling temperature (0.0--2.0) \\
\code{max\_tokens} & \type{int} & No & Max tokens per response \\
\code{max\_steps} & \type{int} & No & Max tool-calling iterations \\
\code{memory} & memory type & No & Memory configuration (default: \code{none}) \\
\code{handoffs} & \type{[agent]} & No & Agents this agent can delegate to \\
\code{strategy} & \type{string} & No & Handoff strategy: \code{"parallel"} or \code{"sequential"} \\
\code{scope} & \type{string} & No & Text description of what agent is allowed to help with \\
\code{scope\_deny} & \type{string} & No & Rejection message for out-of-scope requests \\
\code{timeout} & \type{int} & No & Timeout in seconds (default: 120) \\
\bottomrule
\end{tabular}
\caption{Agent fields}
\end{table}

\subsection{Multiline System Prompts}

For longer system prompts, use triple-quoted strings:

\begin{lstlisting}
agent SupportBot {
    provider: anthropic
    system: """
        You are a customer support agent for Acme Corp.
        Use the knowledge base to answer questions.
        Look up orders when customers ask about purchases.
        Create tickets for issues you cannot resolve.
        Always be polite and helpful.
    """
    tools: [search_kb, lookup_order, create_ticket]
    temperature: 0.3
}
\end{lstlisting}

\section{Agent Methods}

Agents have three methods for interacting with the LLM:

\subsection{.ask() --- Text In, Text Out}

The simplest interaction. Send a text prompt, get a text response:

\begin{lstlisting}
answer, err = Researcher.ask("What is quantum computing?")
if err != nil {
    io.println("Error: " + err.message)
    return
}
io.println(answer)
\end{lstlisting}

\textbf{Signature:} \code{.ask(prompt: string, session: string?) -> (string, Error?)}

\subsection{.run() --- Structured In, Structured Out}

For typed interactions. Send a struct, get a typed response. The output type is inferred from the left-side type annotation:

\begin{lstlisting}
struct ResearchRequest {
    topic: string
    depth: string = "detailed"
}

struct ResearchResult {
    summary: string
    sources: [string]
    confidence: float
}

research: ResearchResult, err = Researcher.run(ResearchRequest{
    topic: "quantum computing"
})
\end{lstlisting}

The compiler generates JSON schemas for both the input and output types. The agent is instructed to return structured data matching the output schema.

\textbf{Signature:} \code{.run(input: T, session: string?) -> (U, Error?)}

Where \code{U} is inferred from the left-side type annotation.

\subsection{.stream() --- Text In, Streaming Out}

For real-time token-by-token responses. Returns an iterable stream:

\begin{lstlisting}
for chunk in Assistant.stream("Tell me a story") {
    io.print(chunk)
}
io.println()  // newline after stream completes
\end{lstlisting}

\textbf{Signature:} \code{.stream(prompt: string, session: string?) -> stream<string>}

The \code{stream<string>} type is iterable with \keyword{for}. Each iteration yields the next token (or chunk) as the LLM generates it.

\section{Memory}

By default, agents are stateless --- each call is independent. Adding memory enables agents to remember conversation history across calls.

\subsection{No Memory (Default)}

\begin{lstlisting}
agent Classifier {
    provider: local
    system: "Classify inputs into categories."
}

// Each call is independent
cat1 = Classifier.ask("This is a bug report")     // "technical"
cat2 = Classifier.ask("What did I just send you?") // has no context
\end{lstlisting}

\subsection{Conversation Memory}

Keeps the last N turns of chat history:

\begin{lstlisting}
agent Assistant {
    provider: anthropic
    system: "You are a helpful assistant."
    memory: conversation(max_turns: 50)
}
\end{lstlisting}

\subsection{Summary Memory}

Summarizes older messages to stay within a token budget. Useful for long-running conversations:

\begin{lstlisting}
agent ProjectHelper {
    provider: anthropic
    system: "You help with long-running software projects."
    memory: conversation(max_turns: 50)
}
\end{lstlisting}

\subsection{Memory Types}

\begin{table}[h]
\centering
\begin{tabular}{lll}
\toprule
\textbf{Type} & \textbf{Parameter} & \textbf{Behavior} \\
\midrule
\code{none} & --- & Stateless (default) \\
\code{conversation(max\_turns: N)} & N: int & Keeps last N message pairs \\
\bottomrule
\end{tabular}
\caption{Built-in memory types}
\end{table}

\section{Sessions}

When an agent has memory, calls are scoped to a \textbf{session}. A session groups related interactions --- for example, one user's conversation.

\begin{lstlisting}
// With session --- memory persists within same session
answer1, _ = Assistant.ask("My name is Alice", session: "user-123")
answer2, _ = Assistant.ask("What is my name?", session: "user-123")
// answer2 = "Your name is Alice"

// Different session --- separate memory
answer3, _ = Assistant.ask("What is my name?", session: "user-456")
// answer3 = "I don't know your name yet"
\end{lstlisting}

The \code{session:} parameter is a named argument available on all agent methods:

\begin{lstlisting}
// .ask() with session
reply, err = Agent.ask("Hello", session: session_id)

// .run() with session
result: T, err = Agent.run(input, session: session_id)

// .stream() with session
for chunk in Agent.stream("Hello", session: session_id) {
    io.print(chunk)
}
\end{lstlisting}

Without a \code{session:} parameter, each call starts a fresh conversation --- even if the agent has memory configured.

\section{Handoffs}

Agents can transfer a conversation to another agent. This is useful when a generalist agent detects that a specialist should take over.

\subsection{Declaring Handoffs}

Use the \code{handoffs} field to list agents that this agent can delegate to:

\begin{lstlisting}
agent FrontDesk {
    provider: anthropic
    system: """
        You are a front desk agent. Greet users and help them.
        If they have a billing question, hand off to BillingAgent.
        If they have a technical issue, hand off to TechAgent.
    """
    handoffs: [BillingAgent, TechAgent]
    memory: conversation(max_turns: 10)
}

agent BillingAgent {
    provider: anthropic
    system: "You handle billing questions. You can look up invoices and process refunds."
    tools: [lookup_invoice, process_refund]
    memory: conversation(max_turns: 30)
}

agent TechAgent {
    provider: anthropic
    system: "You handle technical support. You can search docs and check system status."
    tools: [search_docs, check_status]
    memory: conversation(max_turns: 30)
}
\end{lstlisting}

\subsection{How Handoffs Work}

When an agent has \code{handoffs}, the runtime injects synthetic tools for each target agent (e.g., \code{transfer\_to\_BillingAgent}). The LLM decides when to call these tools based on the system prompt and conversation context.

The handoff process:

\begin{enumerate}
    \item The LLM calls a handoff tool (e.g., \code{transfer\_to\_BillingAgent})
    \item The runtime transfers the conversation history to the target agent's session
    \item The target agent processes the original message with full context
    \item The target agent's response is returned to the caller
\end{enumerate}

\subsection{.ask() --- Automatic Handoffs}

When using \code{.ask()}, handoffs are automatic. The caller sees a seamless response regardless of which agent ultimately handled it:

\begin{lstlisting}
@webhook("/api/chat")
workflow Chat(message: string, session_id: string) -> { reply: string } {
    // FrontDesk may hand off to BillingAgent or TechAgent internally
    // The caller just gets the final reply
    reply, err = FrontDesk.ask(message, session: session_id)
    if err != nil { return { reply: "Something went wrong." } }
    return { reply: reply }
}
\end{lstlisting}

\subsection{.run() --- Manual Handoff Control}

When using \code{.run()}, the workflow receives an \code{AgentResult} that includes handoff information, allowing manual inspection or override:

\begin{lstlisting}
@webhook("/api/chat")
workflow Chat(message: string, session_id: string) -> { reply: string } {
    result: AgentResult, err = FrontDesk.run(message, session: session_id)
    if err != nil { return { reply: "Something went wrong." } }

    if result.handed_off_to != nil {
        io.println("Handed off to: " + result.handed_off_to)
    }

    return { reply: result.reply }
}
\end{lstlisting}

\subsection{Handoff Chains}

Target agents can themselves have handoffs, creating delegation chains. The runtime follows the chain until an agent produces a final response (up to a configurable depth limit):

\begin{lstlisting}
agent L1Support {
    provider: anthropic
    system: "You handle basic support. Escalate complex issues to L2."
    handoffs: [L2Support]
}

agent L2Support {
    provider: anthropic
    system: "You handle complex technical issues. Escalate critical outages to L3."
    tools: [search_docs, check_status]
    handoffs: [L3Support]
}

agent L3Support {
    provider: anthropic
    system: "You handle critical outages. You have full system access."
    tools: [search_docs, check_status, restart_service, page_oncall]
}
\end{lstlisting}

\section{Tool-Calling Loop}

When an agent has tools, the runtime executes a loop:

\begin{enumerate}
    \item Send the prompt (and conversation history) to the LLM
    \item If the LLM requests a tool call, execute the tool
    \item Send the tool result back to the LLM
    \item Repeat until the LLM produces a final response (or \code{max\_steps} is reached)
\end{enumerate}

This is transparent to the caller:

\begin{lstlisting}
// The agent may call search() multiple times internally
// before producing a final answer
answer, err = Researcher.ask("Find recent papers on quantum error correction")
\end{lstlisting}

The \code{max\_steps} field limits how many tool-calling iterations are allowed:

\begin{lstlisting}
agent Researcher {
    provider: anthropic
    system: "You are a thorough researcher."
    tools: [search, read_url]
    max_steps: 10      // at most 10 tool calls per request
    temperature: 0.2
}
\end{lstlisting}

\section{Error Handling}

Agent calls can fail. Always check for errors:

\begin{lstlisting}
answer, err = Researcher.ask("Complex query here")
if err != nil {
    io.println("Agent error: " + err.message)
    // Handle: retry, fallback, or report
}
\end{lstlisting}

Common error cases:
\begin{itemize}
    \item LLM provider unreachable or returns an error
    \item Max steps exceeded (tool-calling loop did not converge)
    \item Structured output did not match expected schema
    \item Timeout exceeded
\end{itemize}

\section{Examples}

\subsection{Stateless Classifier}

\begin{lstlisting}
agent Classifier {
    provider: local
    system: "Classify the input as: positive, negative, or neutral. Reply with just the label."
    temperature: 0.0
}

fn main() {
    label, _ = Classifier.ask("I love this product!")
    io.println(label)  // "positive"
}
\end{lstlisting}

\subsection{Research Agent with Tools}

\begin{lstlisting}
agent Researcher {
    provider: anthropic
    system: "You are a thorough researcher. Always cite your sources."
    tools: [search, read_url]
    temperature: 0.2
    max_steps: 15
}

fn main() {
    answer, err = Researcher.ask("What are the latest advances in fusion energy?")
    if err != nil {
        io.println("Error: " + err.message)
        return
    }
    io.println(answer)
}
\end{lstlisting}

\subsection{Conversational Assistant}

\begin{lstlisting}
agent Tutor {
    provider: anthropic
    system: "You are a patient math tutor. Explain concepts step by step."
    memory: conversation(max_turns: 30)
    temperature: 0.5
}

fn main() {
    session = "student-001"

    reply1, _ = Tutor.ask("What is a derivative?", session: session)
    io.println(reply1)

    reply2, _ = Tutor.ask("Can you give me an example?", session: session)
    io.println(reply2)

    reply3, _ = Tutor.ask("Now explain the chain rule", session: session)
    io.println(reply3)
    // Tutor remembers the full conversation context
}
\end{lstlisting}

\section{Delegation Strategy}

When an agent has \code{handoffs}, the default behavior is LLM-driven: the agent decides when to delegate via synthetic tool calls. The \code{strategy} field overrides this with deterministic routing:

\subsection{Parallel Strategy}

Fan out to all handoff targets concurrently, aggregate their replies:

\begin{lstlisting}
agent ResearchTeam {
    provider: anthropic
    system: "You coordinate research across multiple specialists."
    handoffs: [WebSearcher, PaperAnalyzer, NewsScanner]
    strategy: "parallel"
}

// All three agents run concurrently, results are combined
result, err = ResearchTeam.ask("Latest advances in quantum computing")
\end{lstlisting}

\subsection{Sequential Strategy}

Chain through handoff targets in order, each output feeds the next:

\begin{lstlisting}
agent Pipeline {
    provider: anthropic
    system: "You orchestrate a content pipeline."
    handoffs: [Researcher, Writer, Editor]
    strategy: "sequential"
}

// Researcher -> Writer -> Editor, each builds on the previous
result, err = Pipeline.ask("Write an article about AI safety")
\end{lstlisting}

\begin{table}[h]
\centering
\begin{tabular}{lll}
\toprule
\textbf{Strategy} & \textbf{Behavior} & \textbf{Use Case} \\
\midrule
(default) & LLM decides via tool calls & Dynamic routing \\
\code{"parallel"} & Fan-out to all targets concurrently & Research, voting \\
\code{"sequential"} & Chain targets in order & Pipelines, refinement \\
\bottomrule
\end{tabular}
\caption{Handoff strategies}
\end{table}

\section{Pre-built Agent Templates}

The \code{agents} stdlib package provides reusable agent configurations:

\begin{lstlisting}
import "agents"

// Get a pre-configured template
config = agents.code_reviewer()

// Use with create_agent() to instantiate
reviewer = create_agent(config, claude, [analyze_code])
result, err = reviewer.ask("Review this pull request")
\end{lstlisting}

Available templates: \code{code\_reviewer()}, \code{planner()}, \code{security\_reviewer()}, \code{summarizer()}, \code{data\_analyst()}, \code{customer\_support()}, \code{tdd\_guide()}, \code{doc\_writer()}.

\section{Agent Grammar}

\begin{lstlisting}[language=EBNF]
agent_decl    = "agent" identifier "{" { agent_field } "}" ;

agent_field   = identifier ":" agent_value ;

agent_value   = expression
              | string_lit
              | triple_string
              | "[" ident_list "]"
              | memory_value ;

memory_value  = "conversation" "(" [ arguments ] ")"
              | "summary" "(" [ arguments ] ")"
              | "none" ;
\end{lstlisting}
